\section{Introduction}\label{sec:intro}

Let $(X,Y)$ be a random pair with $Y>0$ and $X\in\R^p$, and suppose that,
conditionally on $X=x$,
\begin{equation}\label{eq:intro-rv}
 \frac{\Pp(Y>ty\mid X=x)}{\Pp(Y>y\mid X=x)}
 \longrightarrow t^{-1/\gamma(x)},\qquad y\to\infty,\quad t>0,
\end{equation}
for a positive function $\gamma$, the conditional tail index. Estimating
$\gamma$ is the first step toward extreme conditional quantiles and tail
probabilities \citep{deHaanFerreira2006,hill1975}, and it suffers from a
particularly severe curse of dimensionality: local estimation must find
observations near $x$ and then keep only an upper fraction of their
responses. A one-dimensional window of width $h$ combined with a tail
fraction $\alpha$ leaves an effective sample of order $n\alpha h$, not $n$
\citep{daouiaEtAl2013,gardesStupfler2014}. When $p$ is large, even moderate
sample sizes are effectively small.

This paper asks the screening question: among $p$ candidate covariates,
possibly with $p$ growing with $n$, which ones does $\gamma$ actually depend
on? Sure independence screening \citep{fanLv2008} answers such questions by
ranking coordinates through inexpensive marginal statistics and keeping a
moderate number for subsequent joint analysis. Model-free variants exist for
general dependence \citep{liZhongZhu2012} and for fixed conditional quantiles
\citep{heWangHong2013}, but the limiting exponent in \eqref{eq:intro-rv} is a
different target: a covariate can move the $0.95$ quantile of $Y$ without
affecting $\gamma$, and conversely. Our simulations make this distinction
concrete.

The screen we propose rests on a single geometric mechanism. Transform each
continuous covariate to $U_j=F_j(X_j)$ and let $\xi_j(u)$ denote the tail
index of $Y$ given $U_j=u$ alone. Conditioning on one coordinate leaves the
others free to range over the fibre $\{u\}\times[0,1]^{p-1}$, so the
projected distribution is a mixture, and a mixture of heavy tails is governed
by its heaviest component. Under a full-support condition,
\begin{equation}\label{eq:intro-envelope}
 \xi_j(u)=\max_{v\in[0,1]^{p-1}}\gamma(u,v):
\end{equation}
the projected index is the upper envelope of the tail-index surface along the
fibre. If coordinate $j$ is inactive, every fibre reaches the global maximum
$\gmax=\max_x\gamma(x)$ and the envelope is flat at $\gmax$. If coordinate
$j$ is active and some of its fibres miss the global argmax set, the envelope
dips below $\gmax$ there. Averaging the envelope over a trimmed interval
gives a score $\Psi_j$ whose gap $\Delta_j=\gmax-\Psi_j$ is zero for every
inactive coordinate and positive for every detectable active one. Ranking
coordinates by increasing $\Psi_j$ is the screen.

The same geometry states the method's limitation. An active coordinate is
invisible when every one of its fibres still meets the global argmax set;
this happens, for instance, for $\gamma(u_1,u_2)=\gamma_0-(u_1-u_2)^2$,
whose argmax is the diagonal. Detectability is a property of the argmax
projections, and we characterize it exactly rather than assuming a generic
signal-strength condition.

Estimation follows the envelope. For each coordinate we slide a window over
the empirical marginal ranks, apply a Hill statistic to the upper
$\alpha$-fraction of the responses in the window, and average the resulting
profile. Empirical ranks give every coordinate the same deterministic window
geometry and remove the marginal scales. The price is a rank-displacement
term, which we control simultaneously over $p$ coordinates by a
Dvoretzky--Kiefer--Wolfowitz bound; it is negligible when $\log(pn)=o(nh^2)$.
The main result bounds the maximal score error by
\[
 O_{\Pp}\!\left(\sqrt{\frac{\log(pn)}{n\alpha h}}\right)+o(\Delta_{\min}),
\]
where $\Delta_{\min}$ is the weakest population gap, and yields sure
screening when $\log(pn)=o(n\alpha h\Delta_{\min}^2)$: ultra-high dimension
is affordable exactly when the local extreme count dominates the logarithm of
the dimension. A blockwise implementation costs $O(pn\log n)$ time.

Our construction adapts the projected-tail framework of
\citet{gardesPodgorny2025}, who estimate a low-dimensional central tail-index
subspace by minimizing an average of projected local Hill estimates over
projection matrices, in fixed dimension. We keep their projection semantics
and local Hill estimator, and change what is needed for screening: the
projections are the $p$ coordinates themselves, $p$ may diverge, windows are
built from empirical ranks with deterministic local counts, and uniformity
over a compact matrix class is replaced by union bounds over coordinates and
window states. Their density assumptions on the covariate are not needed
here; the empirical-rank construction replaces them. Among screening
proposals, \citet{yoshidaUmezu2026} screen on covariate-dependent
extreme-value indices through kernel conditional Pickands contrasts in a
large-bandwidth regime, and parametric tail-index regression \citep{wangTsai2009}, extended to
high dimensions with regularization by \citet{sasakiTaoWang2026},
exploits a parametric link when one is available. Tail dimension reduction methods
\citep{gardes2018,aghbalouEtAl2024,bousebataEtAl2023,arbelEtAl2024} target
related but distinct conditional-tail objects.

The theoretical scope deserves a clear statement. The primitive model
\eqref{eq:intro-rv} identifies the projected index
\eqref{eq:intro-envelope}, but does not by itself provide the uniform
projected-quantile regularity needed for growing-dimensional concentration:
projection can create additional logarithmic slow variation even from exact
Pareto conditionals. We therefore state the required conditions directly on
the projected quantiles, establish them on an explicit model class, and leave
their derivation from primitive assumptions open.

Section~\ref{sec:population} develops the population theory.
Section~\ref{sec:estimation} defines the estimator and states the screening
results. Section~\ref{sec:simulation} reports a simulation study built
around four models whose slowly varying factors are driven by inactive
covariates, a tuning study, one-factor sensitivity analyses, and a
three-method comparison. Section~\ref{sec:realdata} applies the screen to
violent-crime rates in U.S. communities. Section~\ref{sec:discussion}
discusses limitations. Proofs are in Appendix~\ref{app:proofs}.
